\begin{abstract}
Vision--language models for fine-grained understanding must connect image
appearance to explicit, measurable geometry and domain-specific language.
Most benchmarks pair images only with semantic labels or captions and omit
the geometry that grounds them. Grounding between language and geometry
therefore remains untrained and unevaluated. Yoga pose understanding poses
this problem in a demanding form. Asanas differ in subtle joint flexion,
limb elevation, trunk orientation, and bilateral symmetry. Yoga-82, the
standard yoga-pose-recognition benchmark, provides only category labels,
without 3D body measurements or language supervision grounded in them. We
introduce \textbf{YogaAlign}, a benchmark derived from Yoga-82 that augments
yoga images with SAM 3D Body reconstructions, interpretable joint- and
trunk-level descriptors, bilateral-symmetry measures, and instruction-style
annotations. YogaAlign contains 10,090 records, comprising 6,000
image-specific geometry-grounded question--answer pairs and 4,090
pose-level domain-knowledge pairs.

We further propose \textbf{YogaLLaVA V3}, a geometry-conditioned
vision--language model that represents structured measurements as metric-aware
tokens and couples autoregressive language generation with continuous
angle-slot regression. On 521 held-out angle slots, YogaLLaVA V3 obtains a
teacher-forced MAE of $13.98^\circ$, an RMSE of $21.82^\circ$, and 57.01\%
accuracy within $10^\circ$ of the reconstruction-derived references. Under
free-running generation, it emits the correct number of numeric slots for
86.1\% of 310 geometry-grounded test records and achieves an MAE of
$17.65^\circ$ on the 160-record common-success subset shared by all evaluated
models. Because nearly all evaluated targets are supplied through the
structured geometry input, these experiments assess question-conditioned
measurement selection and expression rather than independent biomechanical
angle recovery from RGB images.
\end{abstract}

\begin{IEEEkeywords}
vision--language models; yoga pose understanding; 3D human mesh recovery;
geometry grounding; instruction tuning; Yoga-82
\end{IEEEkeywords}
